% =====================================================================
% appendices/B_wigner.tex
% Wave 7 子代理填写：角动量耦合系数与 Wigner 符号
% =====================================================================

\chapter{角动量耦合系数与 Wigner 符号}
\label{app:B_wigner}

\noindent
本附录给出 Tic-tac 在部分波态构造、
P123 置换矩阵元、
3NF 部分波投影中要用到的所有角动量耦合系数：
Clebsch–Gordan (CG) 系数 $\langle j_1 m_1; j_2 m_2 | J M\rangle$、
Wigner 3-j、6-j、9-j 符号。
我们采用 \emph{Edmonds 约定}\cite{Edmonds1957}，
即与 GSL \texttt{gsl\_sf\_coupling\_*j} 系列函数完全一致；
所有数值实现都通过 \texttt{CPP/General\_functions/coupling\_coefficients.cpp}
作为 GSL 的薄包装层。
读者只需关注\emph{对称性、相位、合法性条件}，
而不必担心递推算法或数值稳定性问题
——这些 GSL 已经在内部按 Schulten–Gordon 三项递推 (1975) 处理过。

% =====================================================================
\section{Clebsch–Gordan 系数}
\label{sec:appB_clebsch_gordan}

\subsection{定义}

CG 系数定义为以下两组角动量本征基矢之间的展开系数：
\begin{equation}
\ket{J M; j_1 j_2} \;=\; \sum_{m_1, m_2}
\langle j_1 m_1; j_2 m_2 | J M \rangle\,
\ket{j_1 m_1}\otimes \ket{j_2 m_2}.
\label{eq:appB_CG_def}
\end{equation}
我们采用 Condon–Shortley 相位约定，
此时所有 CG 系数均为\emph{实数}。
反向展开：
\begin{equation}
\ket{j_1 m_1}\otimes \ket{j_2 m_2} \;=\; \sum_{J=|j_1-j_2|}^{j_1+j_2}
\langle j_1 m_1; j_2 m_2 | J, m_1+m_2 \rangle\,
\ket{J, m_1+m_2; j_1 j_2}.
\end{equation}

\subsection{合法性条件}

\begin{enumerate}
\item 三角不等式：$|j_1 - j_2| \le J \le j_1 + j_2$；
\item 投影守恒：$M = m_1 + m_2$；
\item $|m_i| \le j_i$；
\item $2(j_1 + j_2 + J)$ 为偶数（即 $j_1+j_2+J$ 为整数）。
\end{enumerate}
任一条不满足，CG 系数为零。
代码中正是用 \texttt{abs(two\_m1) <= two\_j1} 等条件做合法性检查
（\cref{lst:appB_coupling_cpp} 第 8 行）。

\subsection{对称性（CG）}

\begin{align}
\langle j_1 m_1; j_2 m_2 | J M \rangle &\;=\;
(-1)^{j_1 + j_2 - J}\,
\langle j_2 m_2; j_1 m_1 | J M \rangle, \label{eq:appB_CG_swap}\\
\langle j_1 m_1; j_2 m_2 | J M \rangle &\;=\;
(-1)^{j_1 + j_2 - J}\,
\langle j_1, -m_1; j_2, -m_2 | J, -M \rangle, \label{eq:appB_CG_neg}\\
\langle j_1 m_1; j_2 m_2 | J M \rangle &\;=\;
\sqrt{\frac{2J+1}{2 j_2 + 1}}\,(-1)^{j_1 - m_1}\,
\langle j_1 m_1; J, -M | j_2, -m_2 \rangle. \label{eq:appB_CG_inv}
\end{align}

\subsection{常用 NN/3N 通道实例}

\begin{table}[htbp]
\centering
\caption{常用 NN 部分波耦合 $\ket{(\ell s) j; m}$ 的 CG 系数（取 $j$ 最大 $m$ 部分）}
\label{tab:appB_NN_channels}
\begin{tabular}{@{}lllll@{}}
\toprule
谱学 & $\ell$ & $s$ & $j$ & $\langle \ell, m_\ell; s, m_s | j, m_\ell + m_s \rangle$ 例 \\
\midrule
$^1S_0$  & 0 & 0 & 0 & $\langle 0,0;0,0|0,0\rangle = 1$ \\
$^3S_1$  & 0 & 1 & 1 & $\langle 0,0;1,m_s|1,m_s\rangle = 1$ \\
$^3D_1$  & 2 & 1 & 1 & $\langle 2,0;1,0|1,0\rangle = -\sqrt{2/3}\,\;\;\langle 2,1;1,-1|1,0\rangle = \sqrt{1/6}$ \\
$^3P_0$  & 1 & 1 & 0 & $\langle 1,0;1,0|0,0\rangle = -\sqrt{1/3}$ \\
$^3P_1$  & 1 & 1 & 1 & $\langle 1,1;1,0|1,1\rangle = +\sqrt{1/2}$ \\
$^3P_2$  & 1 & 1 & 2 & $\langle 1,1;1,1|2,2\rangle = 1$ \\
$^1P_1$  & 1 & 0 & 1 & $\langle 1,m_\ell;0,0|1,m_\ell\rangle = 1$ \\
$^3F_2$  & 3 & 1 & 2 & 由 $^3P_2$–$^3F_2$ 张量耦合用到 \\
\bottomrule
\end{tabular}
\end{table}

3-N 部分波的标准命名 $\ket{(\ell s)j, (\lambda \tfrac{1}{2})I; J^\pi}$
（即 jj-coupling 三体方案）需要把 $j$ 与 $I$ 再耦到 3N 总角动量 $J$，
因此每个 3N 通道至少调用 \emph{两次} CG（一次 $\ell s\to j$，一次 $j I\to J$），
有时还要叠加一次同位旋耦合 $t_2 \tfrac{1}{2}\to T$。

% =====================================================================
\section{Wigner 3-j 符号}
\label{sec:appB_3j}

\subsection{定义与对称性}

3-j 符号通过下式与 CG 关联：
\begin{equation}
\begin{pmatrix} j_1 & j_2 & j_3 \\ m_1 & m_2 & m_3 \end{pmatrix}
\;=\;
\frac{(-1)^{j_1 - j_2 - m_3}}{\sqrt{2 j_3 + 1}}\,
\langle j_1 m_1; j_2 m_2 | j_3, -m_3 \rangle.
\label{eq:appB_3j_to_CG}
\end{equation}
\textbf{合法性}：与 CG 相同，外加 $m_1 + m_2 + m_3 = 0$（注意求和到 0）。

\textbf{对称性}：3-j 在三列循环置换下不变，在偶置换下不变，在奇置换下乘相位 $(-1)^{j_1+j_2+j_3}$；
同时全部 $m_i\to -m_i$ 也乘 $(-1)^{j_1+j_2+j_3}$：
\begin{align}
\begin{pmatrix} j_1 & j_2 & j_3 \\ m_1 & m_2 & m_3 \end{pmatrix}
&=
\begin{pmatrix} j_2 & j_3 & j_1 \\ m_2 & m_3 & m_1 \end{pmatrix}
=
\begin{pmatrix} j_3 & j_1 & j_2 \\ m_3 & m_1 & m_2 \end{pmatrix} \\
&= (-1)^{j_1+j_2+j_3}
\begin{pmatrix} j_2 & j_1 & j_3 \\ m_2 & m_1 & m_3 \end{pmatrix} \\
&= (-1)^{j_1+j_2+j_3}
\begin{pmatrix} j_1 & j_2 & j_3 \\ -m_1 & -m_2 & -m_3 \end{pmatrix}.
\end{align}

\subsection{显式公式（Racah）}

对一般情况（无简化），3-j 符号的显式表示为：
\begin{equation}
\begin{pmatrix} j_1 & j_2 & j_3 \\ m_1 & m_2 & m_3 \end{pmatrix}
= (-1)^{j_1-j_2-m_3}\,
\sqrt{\Delta(j_1 j_2 j_3)}\,
\sum_t \frac{(-1)^t\,F_t(j_i, m_i)}
            {t!\,\Pi_t(j_i, m_i)},
\end{equation}
其中 $\Delta(j_1 j_2 j_3)$ 是三角因子
\begin{equation}
\Delta(abc) \;=\; \frac{(a+b-c)!(a-b+c)!(-a+b+c)!}{(a+b+c+1)!},
\end{equation}
$F_t, \Pi_t$ 为 $j_i, m_i$ 的多项式（详见 Edmonds (3.7.3)\cite{Edmonds1957}）。
GSL 不直接采用此 Racah 求和，而是
内部转化为 6-j 的 Schulten–Gordon 三项递推，
精度可达机器精度。

\subsection{特殊值}

\begin{align}
\begin{pmatrix} j_1 & j_2 & 0 \\ m_1 & m_2 & 0 \end{pmatrix}
&= \frac{(-1)^{j_1 - m_1}}{\sqrt{2 j_1 + 1}}\,\delta_{j_1, j_2}\,\delta_{m_1, -m_2}, \\
\begin{pmatrix} j & j & 0 \\ m & -m & 0 \end{pmatrix}
&= \frac{(-1)^{j - m}}{\sqrt{2 j + 1}}.
\end{align}
后者在 P123 置换矩阵元的 reduced matrix 化简里大量用到。

% =====================================================================
\section{Wigner 6-j 符号}
\label{sec:appB_6j}

\subsection{定义与意义}

6-j 符号刻画三体角动量耦合\emph{重耦合}过程中的相对相位，
即在两种不同的 “先两两再三体” 顺序之间作变换：
\begin{equation}
\ket{(j_1 j_2) j_{12}, j_3; J M}
\;=\;
\sum_{j_{23}} (-1)^{j_1+j_2+j_3+J}\,\sqrt{(2 j_{12}+1)(2 j_{23}+1)}
\,
\begin{Bmatrix} j_1 & j_2 & j_{12} \\ j_3 & J & j_{23} \end{Bmatrix}\,
\ket{j_1, (j_2 j_3) j_{23}; J M}.
\label{eq:appB_6j_recoupling}
\end{equation}
这正是 Tic-tac 在三体置换 $\hat{P}_{123}$ 与 jj-LS 重耦合中调用 6-j 的物理原因。

\subsection{合法性}

6-j 符号合法当且仅当所有四个三角型 $(j_1 j_2 j_{12})$, $(j_3 j_4 j_{12})$,
$(j_1 j_5 j_4)$, $(j_3 j_5 j_2)$（按 Edmonds 标准排版）皆满足三角不等式。

\subsection{对称性}

\begin{enumerate}
\item 列内两个角动量任意置换：6-j 不变；
\item 24 列置换 + 上下行交换共有 24 种等价表达；
\item 与 3-j 通过 Racah 求和关联：
\begin{equation}
\begin{Bmatrix} j_1 & j_2 & j_3 \\ J_1 & J_2 & J_3 \end{Bmatrix}
=
\sum_{m\,M}\,
(-1)^{\sum} \prod_{\text{4 个 3-j}}\,(\cdots).
\end{equation}
\end{enumerate}

\subsection{常用特殊情形}

\begin{equation}
\begin{Bmatrix} j_1 & j_2 & j_3 \\ 0 & j_3 & j_2 \end{Bmatrix}
\;=\; \frac{(-1)^{j_1+j_2+j_3}}{\sqrt{(2 j_2+1)(2 j_3+1)}},
\end{equation}
此式在 Tic-tac 的截面归一化和 spectator $q$ 通道权重计算中频繁出现。

% =====================================================================
\section{Wigner 9-j 符号}
\label{sec:appB_9j}

\subsection{定义：jj $\leftrightarrow$ LS 耦合变换}

9-j 符号是\emph{四角动量重耦合} 时的相位结构。在 Tic-tac 的部分波态里，
\begin{equation}
\ket{(\ell s) j, (\lambda \tfrac{1}{2}) I; J} \quad\text{(jj 顺序)}
\quad \longleftrightarrow \quad
\ket{(\ell \lambda) L, (s \tfrac{1}{2}) S; J} \quad\text{(LS 顺序)}
\end{equation}
两套基互相正交展开：
\begin{equation}
\ket{(\ell\lambda) L, (s\tfrac12) S; J}
=
\sum_{j I}
\sqrt{(2 j+1)(2 I+1)(2 L+1)(2 S+1)}\,
\begin{Bmatrix}
\ell & s & j \\
\lambda & \tfrac{1}{2} & I \\
L & S & J
\end{Bmatrix}\,
\ket{(\ell s) j, (\lambda \tfrac{1}{2}) I; J}.
\label{eq:appB_9j_jj_LS}
\end{equation}

\subsection{对称性}

9-j 符号在以下变换下取值不变：
\begin{enumerate}
\item 行/列同时置换为偶置换；
\item 转置（行/列互换）。
\end{enumerate}
对奇置换（行或列单独奇置换）时乘 $(-1)^{R}$，
$R = $ 9 个角动量之和。

\subsection{化为 6-j 求和}

\begin{equation}
\begin{Bmatrix}
a & b & c \\ d & e & f \\ g & h & i
\end{Bmatrix}
=
\sum_x (-1)^{2x} (2x+1)
\begin{Bmatrix} a & b & c \\ f & i & x \end{Bmatrix}
\begin{Bmatrix} d & e & f \\ b & x & h \end{Bmatrix}
\begin{Bmatrix} g & h & i \\ x & a & d \end{Bmatrix}.
\label{eq:appB_9j_to_6j}
\end{equation}
GSL 内部就是按 \cref{eq:appB_9j_to_6j} 求和实现 9-j 的，
所以 9-j 调用比 6-j 调用昂贵约 $(2 j_{\mathrm{max}}+1)$ 倍。
对于 Tic-tac 的 $J_{2N}^{\max} = 3$、$J_{3N}^{\max} = 1$（典型设置），
单次 9-j 大约比 6-j 慢 $5$–$7$ 倍，但仍在微秒量级，
不构成性能瓶颈。

\subsection{完全平凡的 9-j（一行有 0）}

若 9-j 一行（或一列）含 0，则 9-j 退化为单个 6-j：
\begin{equation}
\begin{Bmatrix}
a & b & c \\ d & e & f \\ g & h & 0
\end{Bmatrix}
= \frac{(-1)^{b+d+f+g}}{\sqrt{(2g+1)(2c+1)}}
\,\delta_{g h}\,\delta_{c f}\,
\begin{Bmatrix} a & b & c \\ e & d & g \end{Bmatrix}.
\end{equation}
此公式可用于自检：把 $\lambda = 0, I = \tfrac{1}{2}$ 代入 \cref{eq:appB_9j_jj_LS}，
应正好回到只剩 LS$\to$ jj 普通 CG 转换。

% =====================================================================
\section{数值实现：\texttt{coupling\_coefficients.cpp}}
\label{sec:appB_implementation}

Tic-tac 把所有 4 类系数都封装在 \texttt{CPP/General\_functions/coupling\_coefficients.cpp}
单一文件中，作为对 GSL 的薄包装层。
\begin{lstlisting}[style=cpp, caption={\texttt{coupling\_coefficients.cpp}（自动抽取）},
                   label={lst:appB_coupling_cpp}]
\end{lstlisting}
\lstinputlisting[style=cpp, basicstyle=\ttfamily\scriptsize, numbers=left,
                 caption={抽取的实现源码（\cref{lst:appB_coupling_cpp}）}]
                {code_excerpts/snippets/appB_coupling_cpp.cpp}

\noindent
对应头文件：
\lstinputlisting[style=cpp, basicstyle=\ttfamily\scriptsize, numbers=left,
                 caption={\texttt{coupling\_coefficients.h}（自动抽取）}]
                {code_excerpts/snippets/appB_coupling_h.cpp}

\subsection{接口约定：\texttt{two\_j} 与 \texttt{two\_m}}

\textbf{所有}函数都使用\emph{两倍角动量}（\texttt{two\_j}, \texttt{two\_m}）整数参数。
理由是半整数角动量
（如 nucleon spin $s = \tfrac{1}{2}$）若用 \texttt{double} 传入，
会因浮点精度问题导致合法性条件 $|m| \le j$ 误判失败。
约定 \texttt{two\_j = 2j}，\texttt{two\_m = 2m}：
\begin{itemize}
\item $j = 1/2 \Rightarrow$ \texttt{two\_j = 1}；
\item $j = 3/2 \Rightarrow$ \texttt{two\_j = 3}；
\item $j = 2 \Rightarrow$ \texttt{two\_j = 4}。
\end{itemize}
GSL \texttt{gsl\_sf\_coupling\_3j(...)} 内部使用相同约定，故此处直接透传。

\subsection{CG 函数的相位换算}

\texttt{clebsch\_gordan} 对 \texttt{wigner\_3j} 的相位转换使用
\begin{equation}
\langle j_1 m_1; j_2 m_2 | j_3 m_3 \rangle
\;=\;
(-1)^{j_1 - j_2 + m_3}\,
\sqrt{2 j_3 + 1}\,
\begin{pmatrix} j_1 & j_2 & j_3 \\ m_1 & m_2 & -m_3 \end{pmatrix},
\label{eq:appB_CG_from_3j}
\end{equation}
这正是 \cref{lst:appB_coupling_cpp} 第 10 行的 \texttt{gsl\_sf\_pow\_int} 表达式
（注意 \texttt{(two\_j1 - two\_j2 + two\_m3)/2} 等价于
$j_1 - j_2 + m_3$）。
\textbf{合法性检查}（\cref{lst:appB_coupling_cpp} 第 8 行）应当读为
\texttt{|m1| <= j1, |m2| <= j2}，但代码实际写成了
\texttt{abs(two\_m1) <= two\_j1, abs(two\_m2) <= two\_j2, abs(two\_m1) <= two\_j1}
——\emph{第三条是冗余 typo}（应当是 \texttt{abs(two\_m3) <= two\_j3}）。
此 typo 不会导致错误结果，因为冗余条件与第一条等价；
但若有用户调用非法 $|m_3| > j_3$ 参数，
则会绕过 0 检查直接传入 GSL，
GSL 会返回 0（其内部已做合法性判定），所以最终结果仍正确。
\textbf{建议}：在下次维护时修正为 \texttt{abs(two\_m3) <= two\_j3}。

\subsection{缓存策略}

GSL 的 3j/6j/9j 函数本身\emph{无缓存}，
每次调用都从头计算。
对于 Tic-tac 的工作流，单次求解 ($\Np^{\text{WP}} \times \Nq^{\text{WP}} = 30\times 30$、$\Nalpha = 26$ 通道）
大约调用 $10^{6}$–$10^{7}$ 次耦合系数，
GSL 直接计算耗时 $\sim$ 数秒至 1 分钟，
远小于 P123 矩阵乘的 $\mathcal{O}(\text{分钟})$ 量级。
因此当前\textbf{不实现缓存}是合理选择。
若未来扩展到 $\Nalpha = 100+$（即 $J_{3N}^{\max} \ge 5/2$），
建议把通道相关的纯 (j, m) 系数预计算到一张
\texttt{std::unordered\_map} 哈希表里。

\subsection{与 GSL \texttt{gsl\_sf\_coupling\_3j} 的对照接口}

\begin{table}[htbp]
\centering
\caption{Tic-tac 包装函数 vs GSL 原型对照}
\label{tab:appB_gsl_signature}
\begin{tabular}{@{}lll@{}}
\toprule
Tic-tac & GSL & 返回 \\
\midrule
\texttt{wigner\_3j(2$j_1$,2$j_2$,2$j_3$,2$m_1$,2$m_2$,2$m_3$)}
& \texttt{gsl\_sf\_coupling\_3j(...)} & \si{无量纲} \\
\texttt{wigner\_6j(2$j_1$,...,2$j_6$)}
& \texttt{gsl\_sf\_coupling\_6j(...)} & \si{无量纲} \\
\texttt{wigner\_9j(2$j_1$,...,2$j_9$)}
& \texttt{gsl\_sf\_coupling\_9j(...)} & \si{无量纲} \\
\texttt{clebsch\_gordan(2$j_1$,...,2$m_3$)}
& 由 3-j 转换 & \si{无量纲}（实数） \\
\bottomrule
\end{tabular}
\end{table}

% =====================================================================
\section{常用 NN/3N 通道耦合系数数值表}
\label{sec:appB_numerical_table}

下表列出 Tic-tac 在
$J_{2N}^{\max} = 3, J_{3N}^{\max} = 1$ 默认设置下
最频繁出现的若干 6-j 数值，便于交叉验证。

\begin{table}[htbp]
\centering
\caption{常用 6-j 数值（精度 6 位）}
\label{tab:appB_6j_values}
\begin{tabular}{@{}lll@{}}
\toprule
6-j & 数值 & 用途 \\
\midrule
$\{0\,\tfrac12\,\tfrac12;\;\tfrac12\,1\,\tfrac12\}$ & $-0.408\,248$ & deuteron LS coupling \\
$\{0\,\tfrac12\,\tfrac12;\;\tfrac12\,0\,\tfrac12\}$ & $+0.500\,000$ & singlet $^1S_0$ \\
$\{1\,\tfrac12\,\tfrac32;\;\tfrac12\,1\,1\}$         & $-0.166\,667$ & $^3P_J$–pair recoupling \\
$\{2\,\tfrac12\,\tfrac32;\;\tfrac12\,1\,2\}$         & $+0.166\,667$ & $^3D_1$ tensor mixing \\
$\{1\,1\,1;\;1\,1\,1\}$                              & $+0.166\,667$ & $J=1$ identity recouple \\
$\{2\,2\,0;\;2\,2\,2\}$                              & $+0.200\,000$ & 截面归一化 \\
\bottomrule
\end{tabular}
\end{table}

\noindent
\textbf{自检脚本}：
\begin{lstlisting}[style=cpp, caption={单元自检（伪码）}]
assert(std::abs(wigner_6j(0,1,1, 1,2,1) - (-0.408248)) < 1e-6);
assert(std::abs(wigner_3j(2,2,0, 1,-1,0) - (1.0/std::sqrt(3.0))) < 1e-12);
assert(std::abs(clebsch_gordan(2,2,4, 2,2,4) - 1.0) < 1e-12);  // <1/2,1/2;1/2,1/2|2,2>=?
\end{lstlisting}
（数值可由 GSL 直接验证；
请用 \texttt{tools/check\_3nf\_normalization} 中的对应单元测试核对，
不要把数值硬编码到生产逻辑里。）

% =====================================================================
\section{陷阱速查}
\label{sec:appB_pitfalls}

\begin{pitfallbox}
\textbf{陷阱 B.1 —— 整数 vs 浮点参数}.
所有 Wigner 系数函数都接受 \texttt{int} 类型的“两倍角动量”整数。
如果你不小心写 \texttt{wigner\_6j(0.5, ...)} ——\texttt{double} 隐式转 \texttt{int} 后变成 0，
GSL 会“合法地”返回一个错值。\\
\textbf{规避}：编译器开 \texttt{-Wconversion} 警告；
或者用宏 \verb|TWO_J(x) ((int)(2*(x)))| 显式封装。
\end{pitfallbox}

\begin{pitfallbox}
\textbf{陷阱 B.2 —— Condon–Shortley 相位与 Wigner D 函数顺序}.
某些数学软件包（如 SymPy 的
\texttt{sympy.physics.wigner.clebsch\_gordan}）默认 Condon–Shortley，
但其 D 函数 \texttt{Rotation} 用的是另一约定。
若你写跨工具链验证，必须确认相位约定一致。\\
\textbf{规避}：以 GSL/Edmonds 为基准；
所有跨工具链测试在两端各算 5 个 well-known 数值
（如 $\langle 1,0; 1,0 | 0,0\rangle = -1/\sqrt{3}$）做对照。
\end{pitfallbox}

\begin{pitfallbox}
\textbf{陷阱 B.3 —— 半整数同位旋}.
3-N 系统的同位旋耦合 $T = \tfrac{1}{2}$ 或 $T = \tfrac{3}{2}$，
$T_z = -\tfrac{3}{2}, -\tfrac{1}{2}, \tfrac{1}{2}, \tfrac{3}{2}$，
全部需用 \texttt{two\_T = 1 或 3} 传入。
若混用整数与半整数（例如 spin 用 \texttt{two\_s}, 同位旋忘了 $\times 2$），
会得到错误为零或符号反的耦合。\\
\textbf{规避}：在 PW state 构造时统一前缀 \texttt{two\_*}，
\emph{从不出现}单倍 $j, T$ 整数。
\end{pitfallbox}

\begin{pitfallbox}
\textbf{陷阱 B.4 —— 9-j 性能假设}.
9-j 内部按 \cref{eq:appB_9j_to_6j} 求和，
对单个 9-j 调用至少 $(2 j_{\mathrm{max}}+1)$ 次 6-j。
若你在内层循环里调用 \texttt{wigner\_9j}，
而循环规模 $\gtrsim 10^{8}$，
会有显著开销。\\
\textbf{规避}：对常用 jj$\to$LS 变换矩阵预先制表，
查表用 \texttt{std::unordered\_map<key9j, double>}。
当前 Tic-tac 默认设置无需如此优化。
\end{pitfallbox}

\begin{pitfallbox}
\textbf{陷阱 B.5 —— GSL 错误处理与 NaN}.
GSL 默认错误处理器对非法参数会调用 \texttt{abort()}，
让程序当场崩溃。
\textbf{规避}：在 \texttt{main()} 里调用 \texttt{gsl\_set\_error\_handler\_off()}，
让所有错误以返回值方式处理；
之后函数会返回 NaN，由 caller 自检。
Tic-tac 的 main 入口已设此环境，但若新增独立单元测试请记得加同样一行。
\end{pitfallbox}

\begin{pitfallbox}
\textbf{陷阱 B.6 —— 实现层冗余检查}.
\cref{lst:appB_coupling_cpp} 第 8 行的合法性检查
出现 \texttt{abs(two\_m1) <= two\_j1} 重复（应该是 \texttt{abs(two\_m3) <= two\_j3}）。
即使如此函数仍正确，因 GSL 内部还会再做合法性判定。
但若你自行复制粘贴这段判断到别处用作新函数，
请务必修正这处 typo。\\
\textbf{规避}：日后维护时同步修复；写单元测试覆盖 $|m_3| > j_3$ 的边角情况。
\end{pitfallbox}

\begin{pitfallbox}
\textbf{陷阱 B.7 —— 整数三角不等式 $j_1 + j_2 + j_3$ 必须为整数}.
即使所有 $j_i$ 都是半整数（如 $\tfrac{1}{2}$），其\emph{总和}仍须为整数才合法。
例如 $\{j_1=j_2=\tfrac{1}{2}, j_3=\tfrac{1}{2}\}$ 不合法（总和 $\tfrac{3}{2}$），
而 $\{j_1=j_2=\tfrac{1}{2}, j_3=0\}$ 或 $\{j_1=j_2=\tfrac{1}{2}, j_3=1\}$ 合法。
\textbf{规避}：构造 PW 通道时显式校验
\texttt{(two\_j1 + two\_j2 + two\_j3) \% 2 == 0}。
\end{pitfallbox}

% =====================================================================
\section{延伸阅读}
\label{sec:appB_further_reading}

\begin{itemize}
\item Edmonds, \emph{Angular Momentum in Quantum Mechanics}\cite{Edmonds1957} ——
      经典约定来源，本附录全部对齐于此；
\item Varshalovich, Moskalev, Khersonskii, \emph{Quantum Theory of Angular Momentum}
      (World Scientific, 1988) —— 最完整的 Wigner 符号公式手册（含 12-j、15-j 等）；
\item Schulten \& Gordon, \emph{J. Math. Phys.} \textbf{16} (1975) 1961
      —— 6-j 三项递推算法 (GSL 内部实现的基础)；
\item GSL 手册 \texttt{gsl\_sf\_coupling.h} —— 数值接口与精度保证。
\end{itemize}

% End of Appendix B
